\paragraph*{04.05.02.03 Anforderungen und Abnahmekriterien priorisieren und darüber entscheiden}
\label{para:04.05.02.03_AnforderungenUndAbnahmekriterienPriorisierenUndDavonEntscheiden}

% Task
\par Die Anbindung der \gls{sst} an \gls{s4} war im \gls{gls-rup} nicht vorgesehen. Ich definierte, priorisierte und entschied über Anforderungen und Abnahmekriterien für die \gls{sst}-Integration, um die erfolgreiche Umsetzung sicherzustellen, ohne die Ziele des \gls{s4}-Programms zu gefährden.

% Action
\par Im Rahmen einer \gls{cc} ohne \gls{aif} entwickelte ich ein angepasstes Projektdurchführungskonzept, das die \gls{sst}-Integration als festen Projektbestandteil verankerte. Dadurch konnten die Anforderungen und Abnahmekriterien für den \gls{fb} vordefiniert, priorisiert und in das Projektmanagementtool überführt werden (vgl. \autoref{fig:TestpläneSST}).

% Result
\par Mein Dashboard im Projektmanagementtool ermöglichte die kontinuierliche Überwachung der Schnittstellenentwicklung und proaktives Eingreifen bei Planabweichungen. So wurde die \gls{sst}-Integration in Welle 1.0 erfolgreich umgesetzt. Die von mir verantworteten \gls{ricefw}-Objekte blieben im Zeitplan für Welle 1.1 und den anstehenden \gls{it1}, ein wichtiger Meilenstein für die Integration der \gls{sst}s in \gls{s4}.

% Reflect
\par Aufgrund der zweiten Verschiebung des Go-Live-Termins für den \gls{bk} 2400 und der Verschiebung von Welle 1.0 besteht nun die Chance, Anforderungen und Abnahmekriterien für den \gls{fb} weiter zu verfeinern. Damit kann sichergestellt werden, dass die Integration nicht nur technisch erfolgreich umgesetzt wird, sondern auch Prozesswissen in den Betrieb überführt und das \gls{s4}-Programm langfristig entlastet.